\documentclass[a4paper,12pt]{article}
\usepackage[T2A]{fontenc}
\usepackage[utf8]{inputenc}
\usepackage[russian]{babel}
\usepackage{amsmath,amssymb,icomma}
\usepackage[a4paper,left=20mm,right=20mm,top=21mm,bottom=22mm,headheight=25pt,headsep=6mm]{geometry}
\usepackage{microtype,parskip,setspace}
\usepackage{xcolor,tikz,booktabs,tabularx,enumitem,fancyhdr,titlesec}
\usetikzlibrary{arrows.meta,positioning,shapes.geometric,calc}
\definecolor{accent}{HTML}{008080}
\definecolor{darkaccent}{HTML}{004D4D}
\definecolor{textgray}{HTML}{333333}
\color{textgray}
\onehalfspacing
\setlength{\parindent}{0pt}
\setlength{\parskip}{5pt}
\setlist[itemize]{leftmargin=*,itemsep=1pt,topsep=3pt}
\setlist[enumerate]{leftmargin=*,itemsep=8pt,topsep=3pt}
\pagestyle{fancy}
\fancyhf{}
\fancyhead[L]{\small Текстовые задачи: движение}
\fancyhead[R]{\small Ярослав Верещагин}
\fancyfoot[C]{\small\thepage}
\renewcommand{\headrulewidth}{0.45pt}
\renewcommand{\headrule}{\hbox to\headwidth{\color{accent}\leaders\hrule height\headrulewidth\hfill}}
\titleformat{\section}{\Large\color{darkaccent}}{\thesection.}{0.5em}{}[\vspace{2pt}{\color{accent}\titlerule[0.4pt]}]
\titlespacing*{\section}{0pt}{11pt}{8pt}
\newcommand{\topic}[1]{\par\medskip\noindent{\color{darkaccent}\textbf{#1}}\par\smallskip}
\newcommand{\unit}[1]{\,\text{#1}}
\newcommand{\ans}[1]{\textbf{Ответ:} #1.}
\begin{document}
\begin{titlepage}
\thispagestyle{empty}
\centering
\vspace*{30mm}
{\Large\color{darkaccent} МАТЕМАТИКА\par}
\vspace{14mm}
{\Huge Чек-лист\par}
\vspace{11mm}
{\LARGE\color{accent} Текстовые задачи на движение\\[3pt] протяжённых тел и эскалатора\par}
\vspace{16mm}
{\large Подготовка к ОГЭ по математике\par}
\vfill
{\large Ярослав Верещагин\par}
\vspace{9mm}
{26 ноября 2026 г.\par}
\vspace{15mm}
{\small Лёвин Артём Александрович\\эксклюзивно для Ярослава Верещагина\par}
\vspace{17mm}
\begin{tikzpicture}
\draw[accent,line width=.65pt] (0,0) .. controls (2,0.8) and (4,-0.8) .. (6,0) .. controls (8,0.8) and (10,-0.8) .. (12,0);
\fill[accent] (6,0) circle (1.1pt);
\end{tikzpicture}
\vspace*{3mm}
\end{titlepage}
\setcounter{page}{1}

\section{Основные величины и единицы измерения}
При равномерном движении путь $S$, скорость $v$ и время $t$ связаны формулами
\[
 S=vt,\qquad t=\frac{S}{v}\ (v>0),\qquad v=\frac{S}{t}\ (t>0).
\]
В задачах удобно приводить длину и время к единицам, в которых дана скорость. Если скорость выражена в км/ч, подставляйте расстояния в километрах, а время --- в часах. Если скорость выражена в м/с, используйте метры и секунды.
\begin{center}
\begin{tabularx}{\linewidth}{@{}lXl@{}}
\toprule
Исходная единица & Действие & Пример \\
\midrule
Метры $\to$ километры & разделить на $1000$ & $900\unit{м}=0,9\unit{км}$\\
Секунды $\to$ часы & разделить на $3600$ & $30\unit{с}=\frac{30}{3600}\unit{ч}$\\
Минуты $\to$ часы & разделить на $60$ & $14\unit{мин}=\frac{14}{60}\unit{ч}$\\
Километры $\to$ метры & умножить на $1000$ & $0,5\unit{км}=500\unit{м}$\\
\bottomrule
\end{tabularx}
\end{center}
\section{Средняя скорость: короткое повторение}
Средняя скорость равна отношению \emph{всего пройденного пути} ко \emph{всему затраченному времени} (включая остановки, если они входят в рассматриваемый промежуток):
\[
 v_{\text{ср}}=\frac{S_1+S_2+\dots+S_n}{t_1+t_2+\dots+t_n}.
\]
\topic{Пример 1. Равные промежутки времени}
Автомобиль ехал по одному часу со скоростями $55$, $65$ и $90\unit{км/ч}$. За три часа он проехал $55+65+90=210\unit{км}$. Поэтому
\[
 v_{\text{ср}}=\frac{210}{3}=70\unit{км/ч}.
\]
\textbf{Важно:} среднее арифметическое скоростей применимо при равных промежутках времени. В общем случае сначала найдите суммарный путь и суммарное время.

\clearpage
\section{Движение протяжённых тел}
Материальная точка --- тело, размерами которого в условиях данной задачи можно пренебречь. Длину поезда при определении времени его полного прохождения мимо столба или другого поезда учитывать необходимо.

Пусть $l_1$ и $l_2$ --- длины объектов вдоль направления движения, а $v_{\text{отн}}$ --- скорость одного объекта относительно другого. При равномерном движении, если в начале ближайшие концы объектов поравнялись, а в конце объекты полностью разошлись, относительный путь равен сумме длин:
\[
 \boxed{\ t=\frac{l_1+l_2}{v_{\text{отн}}}\ }.
\]
\begin{center}
\begin{tikzpicture}[x=1cm,y=1cm,>=Stealth,font=\small]
\foreach \yy/\lab in {0/Начало,-1.65/Конец}{\node[anchor=east] at (-.3,\yy+.25){\lab};\draw[gray!45] (0,\yy-.12)--(10,\yy-.12);}
\draw[accent,fill=accent!8] (1,0) rectangle (4,.5);
\draw[gray,fill=gray!10] (4,0) rectangle (6,.5);
\draw[accent,fill=accent!8] (6,-1.65) rectangle (9,-1.15);
\draw[gray,fill=gray!10] (4,-1.65) rectangle (6,-1.15);
\node at (2.5,.25){$l_1$};\node at (5,.25){$l_2$};
\node at (7.5,-1.4){$l_1$};\node at (5,-1.4){$l_2$};
\draw[->,accent,thick] (1,.77)--(3,.77) node[midway,above]{$v_{\text{отн}}$};
\draw[<->,accent] (1,-2.05)--(6,-2.05) node[midway,below]{$l_1+l_2$};
\draw[dashed,gray] (1,-.12)--(1,-2.05);
\draw[dashed,gray] (6,-1.65)--(6,-2.05);
\end{tikzpicture}
\end{center}
Схема построена относительно второго объекта: задний конец первого смещается на сумму обеих длин. Если один из объектов --- столб или человек, его длину принимают равной нулю.
\topic{Пример 2. Поезд и неподвижный столб}
Поезд движется со скоростью $60\unit{км/ч}$ и полностью проходит мимо столба за $30$ секунд. Найдём длину поезда $x$ в километрах. Длина и скорость столба принимаются равными нулю:
\[
 \frac{30}{3600}=\frac{x+0}{60+0},\qquad
 x=\frac{30\cdot60}{3600}=0,5\unit{км}.
\]
\ans{500 м}

\clearpage
\section{Относительная скорость: сумма или разность}
Знак в знаменателе зависит от направления движения, а не от того, какой из объектов длиннее.
\begin{center}
\begin{tikzpicture}[>=Stealth,font=\small,x=1cm]
\node[anchor=west] at (0,1.3){Навстречу};
\draw[->,very thick,accent] (0,.8)--(2.2,.8) node[midway,above]{$v_1$};
\draw[<-,very thick,gray] (2.6,.8)--(4.8,.8) node[midway,above]{$v_2$};
\node at (2.4,.15){$v_{\text{отн}}=v_1+v_2$};
\node[anchor=west] at (7,1.3){В одном направлении};
\draw[->,very thick,accent] (7,.8)--(9.2,.8) node[midway,above]{$v_1$};
\draw[->,very thick,gray] (9.6,.8)--(11.8,.8) node[midway,above]{$v_2$};
\node at (9.4,.15){$v_{\text{отн}}=|v_1-v_2|$};
\end{tikzpicture}
\end{center}
Для обгона необходимо, чтобы догоняющий двигался быстрее: $v_1>v_2$. При равных скоростях догнать другой объект с ненулевого расстояния невозможно.
\topic{Пример 3. Поезд обгоняет пешехода}
Поезд движется со скоростью $83\unit{км/ч}$, а пешеход идёт в том же направлении со скоростью $3\unit{км/ч}$. Относительная скорость поезда равна $83-3=80\unit{км/ч}$, а не $86\unit{км/ч}$. Если прохождение мимо пешехода заняло $18$ секунд, длина поезда составляет
\[
 l=(83-3)\cdot\frac{18}{3600}=0,4\unit{км}=400\unit{м}.
\]
\topic{Пример 4. Поезд проходит мимо лесополосы}
Поезд со скоростью $72\unit{км/ч}$ полностью проходит вдоль неподвижной лесополосы длиной $900$ м за одну минуту. Обозначим длину поезда через $x$ (км). От момента, когда передняя часть поезда достигла начала лесополосы, до момента, когда задняя часть миновала её конец, он проходит $x+0,9$ км:
\[
 \frac{1}{60}=\frac{x+0,9}{72},\qquad
 x+0,9=1,2,\qquad x=0,3\unit{км}.
\]
\ans{300 м}

\clearpage
\section{Внимательно выбирайте объект наблюдения}
В задаче о двух поездах принципиально различаются два события: один поезд проходит \emph{мимо человека} в другом поезде или \emph{мимо всего другого поезда}. Скорость пассажира относительно земли равна скорости его поезда, если пассажир неподвижен в вагоне.
\topic{Пример 5. Наблюдатель в купе}
Пассажирский поезд едет со скоростью $65\unit{км/ч}$. Навстречу по параллельному пути движется товарный состав со скоростью $25\unit{км/ч}$. Пассажир у окна заметил, что весь товарный состав прошёл мимо него за $10$ секунд. Требуется найти длину товарного состава.
\begin{center}
\begin{tikzpicture}[>=Stealth,font=\small]
\draw[gray,fill=gray!8] (0,0) rectangle (8,.65);
\foreach \x in {1.6,3.2,4.8,6.4}{\draw[gray] (\x,0)--(\x,.65);}
\fill[accent] (3.4,.325) circle (2.5pt);
\node[above] at (3.4,.66){наблюдатель};
\draw[->,gray,thick] (8.2,.325)--(10,.325) node[midway,above]{$65\unit{км/ч}$};
\draw[accent,fill=accent!8] (3,-1.25) rectangle (7,-.6);
\foreach \x in {4,5,6}{\draw[accent] (\x,-1.25)--(\x,-.6);}
\draw[<-,accent,thick] (1.2,-.92)--(2.8,-.92) node[midway,above]{$25\unit{км/ч}$};
\draw[<->,accent] (3,-1.55)--(7,-1.55) node[midway,below]{длина $x$};
\end{tikzpicture}
\end{center}
Время засечено относительно человека, поэтому длина пассажирского поезда в числителе не появляется. Скорости складываются, так как поезда движутся навстречу друг другу:
\[
 \frac{10}{3600}=\frac{x}{65+25},\qquad
 x=\frac{10\cdot90}{3600}=0,25\unit{км}.
\]
\ans{250 м}
\topic{Пример 6. Один поезд полностью обгоняет другой}
Пассажирский и товарный поезда движутся в одном направлении со скоростями $75$ и $45\unit{км/ч}$ соответственно. Длина товарного поезда $800$ м. Полный обгон занимает $2$ минуты. Пусть $x$ --- длина пассажирского поезда в километрах. Тогда
\[
 \frac{2}{60}=\frac{x+0,8}{75-45},\qquad
 x+0,8=1,\qquad x=0,2\unit{км}.
\]
\ans{200 м}

\clearpage
\section{Обгон с начальным и конечным зазорами}
Если до обгона между объектами уже есть промежуток $d_{\text{нач}}$, а после обгона образуется промежуток $d_{\text{кон}}$, их необходимо добавить к сумме длин:
\[
 t=\frac{l_1+l_2+d_{\text{нач}}+d_{\text{кон}}}{v_1-v_2},\qquad v_1>v_2.
\]
Здесь $v_1$ --- скорость догоняющего объекта, $v_2$ --- скорость обгоняемого. Зазоры измеряются между ближайшими концами объектов.
\begin{center}
\begin{tikzpicture}[>=Stealth,font=\small,x=.93cm]
\node[anchor=east] at (-.2,.22){До};
\node[anchor=east] at (-.2,-1.62){После};
\draw[accent,fill=accent!8] (0,0) rectangle (1.5,.45);
\draw[gray,fill=gray!12] (4,0) rectangle (7,.45);
\draw[gray,fill=gray!12] (4,-1.85) rectangle (7,-1.4);
\draw[accent,fill=accent!8] (10,-1.85) rectangle (11.5,-1.4);
\node at (.75,.225){$l_1$};\node at (5.5,.225){$l_2$};
\node at (5.5,-1.625){$l_2$};\node at (10.75,-1.625){$l_1$};
\draw[->,accent,thick] (0,.72)--(1.5,.72);
\draw[<->] (1.5,-.35)--(4,-.35) node[midway,below]{$d_{\text{нач}}$};
\draw[<->] (7,-2.2)--(10,-2.2) node[midway,below]{$d_{\text{кон}}$};
\draw[dashed,gray] (4,.0)--(4,-1.4);
\draw[dashed,gray] (7,.0)--(7,-1.4);
\end{tikzpicture}
\end{center}
\topic{Пример 7. Две баржи}
По морю в одном направлении движутся баржи длиной $70$ и $30$ м. Сначала короткая баржа отстаёт от длинной на $250$ м (от носа короткой до кормы длинной). Через $14$ минут короткая баржа уже опережает длинную на $350$ м (от кормы короткой до носа длинной). На сколько километров в час скорость короткой баржи больше скорости длинной?

Обозначим искомую разность скоростей через $\Delta v$. Переводим расстояния в километры, а время --- в часы:
\[
 \frac{14}{60}=\frac{0,03+0,07+0,25+0,35}{\Delta v}
 =\frac{0,7}{\Delta v}.
\]
Перемножаем крест-накрест:
\[
 14\Delta v=60\cdot0,7,\qquad
 \Delta v=\frac{42}{14}=3\unit{км/ч}.
\]
\ans{3 км/ч}
\textbf{Главный приём:} если требуется найти, на сколько одна скорость больше другой, достаточно определить их разность. Искать каждую скорость отдельно не нужно.

\clearpage
\section{Движение по эскалатору}
Если человек шагает по эскалатору в направлении его движения, его скорость относительно здания равна сумме скорости ходьбы относительно ступеней и скорости эскалатора. Если человек стоит на ступенях, его скорость относительно здания равна скорости эскалатора.
\topic{Пример 8. Три способа спуска}
Лида спускается по движущемуся вниз эскалатору за $24$ секунды. По неподвижному эскалатору с той же скоростью ходьбы относительно ступеней она спускается за $42$ секунды. За сколько секунд Лида спустится, стоя на ступенях движущегося эскалатора?

Пусть $v_\text{л}$ --- скорость ходьбы Лиды относительно ступеней, $v_\text{э}$ --- скорость эскалатора, $t$ --- искомое время. Обе скорости положительны. Составим таблицу (скорости в м/с, время в секундах, путь в метрах).
\begin{center}
\renewcommand{\arraystretch}{1.2}
\begin{tabularx}{\linewidth}{@{}Xccc@{}}
\toprule
Ситуация & Скорость & Время & Путь \\
\midrule
Лида идёт по движущемуся эскалатору & $v_\text{л}+v_\text{э}$ & $24$ & $24(v_\text{л}+v_\text{э})$\\
Лида идёт по неподвижному эскалатору & $v_\text{л}$ & $42$ & $42v_\text{л}$\\
Лида стоит на движущихся ступенях & $v_\text{э}$ & $t$ & $tv_\text{э}$\\
\bottomrule
\end{tabularx}
\end{center}
Длина эскалатора в трёх случаях одинакова. Поэтому достаточно двух независимых равенств:
\[
 \begin{cases}
 24(v_\text{л}+v_\text{э})=42v_\text{л},\\
 42v_\text{л}=tv_\text{э}.
 \end{cases}
\]
Из первого равенства $24v_\text{э}=18v_\text{л}$, следовательно,
\[
 \frac{v_\text{л}}{v_\text{э}}=\frac{24}{18}=\frac43.
\]
Из второго равенства получаем
\[
 t=42\cdot\frac{v_\text{л}}{v_\text{э}}=42\cdot\frac43=56\unit{с}.
\]
\ans{56 с} Отдельные значения скоростей и длину эскалатора находить не требуется.

\clearpage
\section{Алгоритм: протяжённые тела и эскалатор}
\vspace{5mm}
\begin{center}
\begin{tikzpicture}[>=Stealth, node distance=8mm, font=\small,
 start/.style={ellipse,draw=accent,fill=accent!5,minimum height=9mm,minimum width=42mm,align=center},
 actionbox/.style={rectangle,rounded corners=2pt,draw=accent,minimum height=11mm,text width=65mm,align=center,inner sep=5pt},
 choose/.style={diamond,aspect=2.2,draw=accent,fill=accent!4,align=center,text width=35mm,inner sep=1pt},
 arrow/.style={->,draw=darkaccent,line width=.65pt}]
\node[start] (a){Прочитайте условие};
\node[actionbox,below=of a] (b){Установите, что именно требуется найти. Обозначьте неизвестное};
\node[choose,below=of b] (c){Эскалатор?};
\node[actionbox,below left=9mm and 3mm of c,text width=60mm] (d){Составьте таблицу $v,t,S$. Приравняйте пути по эскалатору};
\node[actionbox,below right=9mm and 3mm of c,text width=60mm] (e){Выделите оба объекта, их длины, направление движения и зазоры};
\node[actionbox,below=of e,text width=60mm] (f){Выберите относительную скорость: сумма при встрече, разность при обгоне};
\node[actionbox,below=of f,text width=60mm] (g){Согласуйте единицы. Сложите длины и зазоры: $S_{\text{отн}}$. Запишите $t=S_{\text{отн}}/v_{\text{отн}}$};
\node[actionbox,below=12mm of g,xshift=-36mm] (h){Решите уравнение или систему. Найдите именно запрошенную величину};
\node[actionbox,below=of h] (i){Проверьте знак, смысл результата и единицы ответа};
\node[start,below=of i] (j){Запишите ответ};
\draw[arrow] (a)--(b);\draw[arrow] (b)--(c);
\draw[arrow] (c.west)-| node[pos=.25,above]{да} (d.north);
\draw[arrow] (c.east)-| node[pos=.25,above]{нет} (e.north);
\draw[arrow] (e)--(f);\draw[arrow] (f)--(g);
\draw[arrow] (d.south) |- (h.west);
\draw[arrow] (g.south) |- (h.east);
\draw[arrow] (h)--(i);\draw[arrow] (i)--(j);
\end{tikzpicture}
\end{center}

\clearpage
\section{Контроль оформления и типичные ошибки}
Перед сдачей решения проверьте, что каждое число в уравнении понятно из условия или предыдущих вычислений. Не ограничивайтесь одним числовым ответом, если требуется развёрнутое решение.
\begin{itemize}
\item[$\square$] Вы обозначили неизвестное и указали, в каких единицах оно выражено.
\item[$\square$] Вы определили, мимо какого объекта происходит движение: человека, столба, лесополосы или целого поезда.
\item[$\square$] При встречном движении скорости сложены, при обгоне --- из большей вычтена меньшая.
\item[$\square$] Учтены обе длины, а при наличии зазоров --- также начальный и конечный промежутки.
\item[$\square$] Расстояния и время согласованы с единицами скорости. Ответ переведён в требуемую единицу.
\item[$\square$] При движении по эскалатору использовано равенство путей, а не равенство скоростей или времени.
\item[$\square$] Результат положителен и соответствует вопросу задачи.
\end{itemize}
\topic{Как не ошибиться при решении пропорции}
Пусть для двух барж получено уравнение
\[
 \frac{18}{60}=\frac{0,6}{\Delta v},\qquad \Delta v>0.
\]
Не переставляйте числа наугад. Перемножьте крайние и средние члены пропорции:
\[
 18\Delta v=60\cdot0,6,\qquad
 \Delta v=\frac{36}{18}=2\unit{км/ч}.
\]
Проверка: $0,6/2=0,3$ часа, то есть $18$ минут.
\topic{Как объяснить каждое слагаемое}
Если длина одного объекта $380$ м, а искомая длина второго объекта равна $x$ км, запишите $380\unit{м}=0,38\unit{км}$. Когда относительный путь равен $0,8$ км, уравнение будет выглядеть так:
\[
 x+0,38=0,8,\qquad x=0,8-0,38=0,42\unit{км}.
\]
Число $0,38$ не должно появляться в решении без пояснения, почему именно его нужно вычесть.

\clearpage
\section*{Тренировочный блок}
Выполните задания последовательно. Для задач на движение записывайте перевод единиц, поясняйте выбор относительной скорости и составляйте уравнение или таблицу. Ответ давайте в единицах, указанных в вопросе.
\topic{Средний уровень}
\begin{enumerate}
\item Автобус отправился из города на экскурсию. Первый час он ехал по загородной дороге со скоростью $60\unit{км/ч}$, второй час --- по шоссе со скоростью $72\unit{км/ч}$, а третий час --- по скоростной трассе со скоростью $84\unit{км/ч}$. В течение этих трёх часов автобус не останавливался. Найдите среднюю скорость автобуса на всём маршруте. Ответ дайте в километрах в час.
\item Пассажирский поезд движется по прямому участку железной дороги с постоянной скоростью $72\unit{км/ч}$. Наблюдатель на обочине засёк время с момента, когда передняя часть поезда поравнялась с придорожным столбом, до момента, когда мимо столба прошёл последний вагон. Это время составило $25$ секунд. Найдите длину поезда в метрах.
\item Вдоль железнодорожного пути идёт пешеход со скоростью $4\unit{км/ч}$. Позади него в том же направлении движется поезд со скоростью $76\unit{км/ч}$. От момента, когда начало поезда поравнялось с пешеходом, до момента, когда весь поезд прошёл мимо него, прошло $20$ секунд. Какова длина поезда? Ответ дайте в метрах.
\item Поезд движется со скоростью $90\unit{км/ч}$ мимо лесополосы, протянувшейся вдоль железной дороги на $750$ метров. С момента, когда передняя часть поезда достигла начала лесополосы, до момента, когда его последний вагон миновал конец лесополосы, прошло $48$ секунд. Найдите длину поезда в метрах.
\end{enumerate}
\topic{Уровень выше среднего}
\begin{enumerate}[resume]
\item По прямому железнодорожному пути движется поезд со скоростью $72\unit{км/ч}$. По дорожке вдоль путей навстречу поезду идёт человек со скоростью $3,6\unit{км/ч}$. Пешеход заметил, что с момента встречи с передней частью поезда до момента прохождения мимо последнего вагона прошло $15$ секунд. Найдите длину поезда в метрах.
\end{enumerate}

\clearpage
\section*{Тренировочный блок (продолжение)}
\topic{Уровень выше среднего}
\begin{enumerate}[resume]
\item Пассажирский поезд движется со скоростью $54\unit{км/ч}$. По соседнему параллельному пути навстречу ему едет товарный поезд со скоростью $72\unit{км/ч}$. Пассажир, сидящий у окна в купе, заметил, что с момента появления передней части товарного поезда напротив него до прохождения последнего вагона прошло $12$ секунд. Найдите длину товарного поезда в метрах.
\item По двум параллельным железнодорожным путям навстречу друг другу движутся два поезда. Длина первого поезда равна $240$ м, длина второго --- $360$ м. Оба поезда движутся равномерно, каждый со скоростью $54\unit{км/ч}$. Сколько секунд пройдёт от момента встречи передних частей поездов до момента, когда их последние вагоны полностью разойдутся?
\item По реке в одном направлении движутся две баржи длиной $45$ и $75$ метров. Короткая баржа сначала находится позади длинной; расстояние от носа короткой баржи до кормы длинной равно $180$ метрам. Через $12$ минут короткая баржа полностью обогнала длинную и расстояние от её кормы до носа длинной составило $300$ метров. На сколько километров в час скорость короткой баржи больше скорости длинной?
\end{enumerate}
\topic{Сложный уровень}
\begin{enumerate}[resume]
\item По морю параллельными курсами в одном направлении следуют две баржи длиной $30$ и $90$ метров. В некоторый момент короткая баржа находится позади длинной, а расстояние от её носа до кормы длинной баржи равно $200$ метрам. Скорость короткой баржи на $2,4\unit{км/ч}$ больше скорости длинной. Через $15$ минут короткая баржа оказывается впереди длинной. Каково теперь расстояние от кормы короткой баржи до носа длинной? Ответ дайте в метрах.
\item На станции метро пассажир, шагая вниз по движущемуся вниз эскалатору, достигает нижней площадки за $30$ секунд. Когда эскалатор неподвижен, тот же пассажир проходит весь спуск за $50$ секунд, сохраняя скорость ходьбы относительно ступеней. За сколько секунд он спустится на таком же движущемся эскалаторе, если будет неподвижно стоять на ступенях? Скорость эскалатора считается постоянной.
\end{enumerate}

\clearpage
\section*{Ответы}
\begin{table}[h!]
\centering
\caption{Ответы к тренировочному блоку}
\renewcommand{\arraystretch}{1.12}
\begin{tabularx}{.76\linewidth}{@{}lX@{}}
\toprule
№ & Ответ \\
\midrule
1 & $72\unit{км/ч}$\\
2 & $500\unit{м}$\\
3 & $400\unit{м}$\\
4 & $450\unit{м}$\\
5 & $315\unit{м}$\\
6 & $420\unit{м}$\\
7 & $20\unit{с}$\\
8 & $3\unit{км/ч}$\\
9 & $280\unit{м}$\\
10 & $75\unit{с}$\\
\bottomrule
\end{tabularx}
\end{table}
\vspace{4mm}
\begin{center}
{\color{darkaccent}\large Итоговая памятка}\par\smallskip
\begin{tabularx}{\linewidth}{@{}p{.37\linewidth}X@{}}
\toprule
Сигнал в условии & Что делать \\
\midrule
Равные промежутки времени & Средняя скорость равна среднему арифметическому скоростей.\\
Поезд проходит мимо столба & Учитывать только длину поезда.\\
Навстречу / в одном направлении & Использовать сумму / разность скоростей.\\
Мимо пассажира в купе & Учитывать длину проходящего поезда, но скорость пассажира равна скорости его поезда.\\
Начальный и конечный зазоры & Прибавить оба зазора к сумме длин.\\
На сколько скорость больше? & Найти разность скоростей.\\
Эскалатор & Составить таблицу и приравнять пути.\\
\bottomrule
\end{tabularx}
\end{center}
\end{document}
